\lamina{img/il_tp.jpg}
\capitulo{CAPÍTULO 9}{LA TABLA DE LAS EQUIVALENCIAS}{O: La piedra de Rosetta de la conciencia}{}
\noindent \primeras{Para}{ consolidar las intuiciones de la primera parte del libro, es útil trazar un mapa comparativo. Este experimento no propone metáforas aisladas: sugiere que la conciencia comparte una \textbf{topología informacional} común con otras estructuras del universo físico, lógico y espiritual.}

La encapsulación no es una propiedad exclusiva del software, ni la evaporación pertenece solo a la termodinámica de los agujeros negros, ni el vacío es un concepto de uso exclusivo de los místicos orientales. Todos estos lenguajes describen las mismas transiciones de fase y límites de acceso bajo escalas diferentes.

La siguiente tabla actúa como la piedra de Rosetta de este libro, traduciendo los conceptos centrales del horizonte a sus equivalentes exactos en tres dominios fundamentales (tabla en la página siguiente).

\begin{landscape}
\thispagestyle{empty}
\vspace*{\fill}
{\scriptsize\renewcommand{\arraystretch}{1.25}\setlength{\tabcolsep}{3.5pt}
\noindent\begin{tabularx}{\linewidth}{>{\raggedright\arraybackslash\hsize=0.70\hsize}X>{\raggedright\arraybackslash\hsize=1.0\hsize}X>{\raggedright\arraybackslash\hsize=1.0\hsize}X>{\raggedright\arraybackslash\hsize=1.0\hsize}X>{\raggedright\arraybackslash\hsize=1.30\hsize}X}\toprule
\thead{Dimensión de Análisis} & \thead{Filosofía del Horizonte} & \thead{Ingeniería de Software} & \thead{Física de Agujeros Negros} & \thead{Filosofías Orientales} \\ \midrule
\textsc{El Origen / El Vacío} & El Reservorio (potencial puro de información desestructurada). & El \emph{Heap} (memoria dinámica global sin asignar, pool común). & Vacío cuántico / Campo subyacente y fluctuaciones. & \emph{Hun Dun} (caos primitivo, Taoísmo) / \emph{Shunyata} (Vacío, Budismo) / \emph{Brahman nirguna} (Hinduismo). \\[4pt]
\textsc{Nacimiento / Identidad} & Condensación (establecimiento de la frontera circular del yo). & Instanciación y Encapsulamiento (creación del objeto, variables \texttt{private}). & Colapso gravitatorio / Formación de horizonte de sucesos (Schwarzschild). & Los 7 orificios en \emph{Hun Dun} / \emph{De}, la eficacia mínima que sobrevive al colapso (Taoísmo) / \emph{Maya} / \emph{Brahman saguna} (Hinduismo). \\[4pt]
\textsc{Experiencia Subjetiva} & Estado privado (perspectiva de primera persona accesible desde dentro). & Estado interno (\texttt{private fields} ocultos tras la interfaz). & Región interior del horizonte (inaccesible para observador externo). & \emph{Atman} (el sí mismo, Hinduismo) / Haz de percepciones (Hume/Budismo). \\[4pt]
\textsc{Conducta / Apariencia} & Interfaz pública (conducta observable, señales neurales). & Interfaz pública / API pública (\texttt{public methods} de interacción). & Superficie del horizonte (codificación holográfica de la entropía). & \emph{Rupa} (forma material externa, Budismo) / \emph{Samsara} (mundo de apariencias). \\[4pt]
\textsc{Muerte / Liberación} & Evaporación (disolución del límite y retorno de la información). & Recolección de basura (\emph{Garbage Collection} / deallocación). & Evaporación de Hawking (pérdida de masa hasta disolución total). & \emph{Nirvana} (cese de Samsara, Budismo) / \emph{Moksha} (la gota que se disuelve en el océano). \\[4pt]
\textsc{Huellas / Legado} & Radiación del yo / Alteración del campo local residual. & Fragmentación, efectos secundarios y registros de \emph{log}. & Radiación de Hawking / Remanente cuántico de información. & \emph{Karma} (huellas y fragmentación que condicionan futuras instanciaciones). \\[4pt]
\textsc{El Vínculo / Conexión} & Resonancia / Entrelazamiento de horizontes. & Acoplamiento de objetos (Asociación, Agregación, Composición). & Entrelazamiento cuántico / Puente Einstein-Rosen (ER=EPR). & Interexistencia (Budismo) / La red de joyas de Indra. \\[4pt]
\textsc{Sincronización} & Enganche de fase / Sintonía temporal subjetiva. & Llamada síncrona / Sincronía por reloj de referencia (NTP/PTP). & Coherencia de fase en osciladores gravitacionales. & \emph{Spanda} (vibración primordial creadora de resonancia, Shivaismo). \\
\bottomrule\end{tabularx}}
\vspace*{\fill}
\end{landscape}

\seccion{La naturaleza de la asimetría}

\noindent El valor de esta tabla no reside en la curiosidad de sus paralelismos, sino en lo que revela sobre la estructura de la realidad. Si el “problema duro” de la conciencia se siente irresoluble, es porque solemos enfocarlo desde una ontología plana: asumimos que todo lo que existe debe ser públicamente observable.

El modelado de sistemas complejos y la física de horizontes nos muestran que esto no es así. El universo es asimétrico. Diseñar sistemas complejos (sean códigos informáticos, horizontes de sucesos cósmicos o cerebros biológicos) requiere, por necesidad termodinámica y lógica, trazar límites de ocultamiento de información. La subjetividad no es un ingrediente místico inyectado en el cerebro; es la perspectiva inevitable de cualquier dominio que ha logrado cerrarse lo suficiente como para albergar un estado privado.

\begin{notacap}{Nota al capítulo 9}
\lbl{Lo que sí sabemos:} Que los sistemas de información, independientemente de su sustrato (silicio, gravedad o carbono), exhiben propiedades isomórficas de frontera cuando alcanzan ciertos niveles de integración causal.

\lbl{Lo que no sabemos:} Si esta concordancia estructural apunta a una ley informacional unificada que subyace a la física, a la computación y a la mente, o si es el límite de la capacidad explicativa de nuestra propia estructura cerebral.

\lbl{Si solo te quedas con una idea:} El yo no es una sustancia. Es una topología de acceso. No hay un “fantasma en la máquina”; lo que hay es un estado privado protegido por la frontera que la naturaleza y los ingenieros usan para evitar que los sistemas colapsen en el caos.

\lbl{Lecturas:} Chalmers, D. (1995), “Facing up to the problem of consciousness”; Tononi, G. (2015), “Integrated Information Theory: From Consciousness to its Physical Substrate”; Parnas, D. (1972), “On the Criteria to be Used in Decomposing Systems into Modules”.
\end{notacap}

